\chapter{Դերեր և խնդիրներ}
\label{ch:tasks}

\section{\studentlabel{1} — Միկրոէլեկտրոնիկա}
\textbf{Անհատական թեմա.} \emph{AUP ZU3 FPGA-ի հիման վրա սենսորային տվյալների գեներացման,
հավաքագրման և նախնական մշակման ապարատային համակարգի նախագծում.}

\begin{tasklist}{1}
  \task Նախագծել կոնֆիգուրելի RTL ազդանշան-գեներատորներ (Ramp, Periodic)։
  \task Իրականացնել սենսոր-էմուլյացիայի տրամաբանությունը (scaling, offset, saturation)։
  \task Իրականացնել LFSR-ի վրա հիմնված աղմուկի գեներատոր։
  \task Իրականացնել անոմալիաների ներարկում։
  \task Մշակել տվյալների հավաքագրման և packetizer տրամաբանությունը։
  \task Իրականացնել FIFO/BRAM բուֆերացում։
  \task Մշակել 3 կոնֆիգուրելի processing IP։
  \task Իրականացնել AXI4-Lite կոնֆիգուրացիոն ռեգիստրները (Աղյուսակ~\ref{tab:regmap})։
  \task Իրականացնել PL$\rightarrow$PS տվյալների փոխանցումը AXI-ով և ինտեգրել PL-ը Zynq PS-ի հետ։
  \task Գրել PS-ի ծրագիրը (Vitis)՝ տվյալները Ethernet-ով ուղարկելու և հրամաններ ընդունելու համար։
  \task Մշակել testbench-եր յուրաքանչյուր մոդուլի համար։
  \task Կատարել սինթեզ, implementation, timing analysis և LUT/FF/DSP/BRAM օգտագործման վերլուծություն։
  \task Վալիդացնել ամբողջական նախագիծը AUP ZU3-ի վրա։
  \task \optionaltask Իրականացնել FIR ֆիլտր կամ AXI DMA ուղի։
\end{tasklist}

\keypoint{Ուսանող 1-ի դիպլոմը կենտրոնանում է թվային նախագծման, RTL-ի, FPGA ճարտարապետության,
data/control ուղիների, բուֆերացման, AXI-ի, վերիֆիկացիայի, սինթեզի, timing-ի և ապարատային
վալիդացիայի վրա։}

\section{\studentlabel{2} — Ծրագրային ապահովում / Ալգորիթմներ}
\textbf{Անհատական թեմա.} \emph{FPGA-ից ստացվող սենսորային տվյալների պահպանման,
վիզուալիզացիայի և վերլուծության ծրագրային հարթակի մշակում.}

\begin{tasklist}{2}
  \task Նախագծել տվյալների բազայի մոդելը (Experiment, Sensor, Sample)։
  \task Մշակել Backend/API՝ տվյալների ընդունման և կոնֆիգուրացիոն հրամանների համար։
  \task Գրել \textbf{FPGA էմուլյատոր} (սկրիպտ), որը ուղարկում է նույն ֆորմատի տվյալներ, որպեսզի աշխատանքը չկախվի ապարատից։
  \task Մշակել Web հավելված՝ փորձի կոնֆիգուրացիայի ինտերֆեյսով։
  \task Իրականացնել live և պատմական տվյալների վիզուալիզացիա և փորձերի համեմատում։
  \task Իրականացնել վիճակագրական վերլուծություն (mean, variance, std, histogram)։
  \task Իրականացնել անոմալիաների հայտնաբերման մեկ մեթոդ և համեմատել այն FPGA-ի threshold detector-ի հետ։
  \task \optionaltask Իրականացնել կանխատեսում (prediction) կամ FPGA-ի պարամետրերի ավտոմատ փոփոխում։
\end{tasklist}

\keypoint{Ուսանող 2-ի դիպլոմը կենտրոնանում է տվյալների բազայի, Backend-ի, Web հարթակի,
վիզուալիզացիայի, վիճակագրության և վերլուծության/մոդելավորման վրա և անկախ է RTL-ի մշակումից։}

\section{Ընդհանուր խնդիրներ և միջերեսներ}
\begin{tasklist}{0}
  \task Մինչև 5-րդ շաբաթը համաձայնեցնել և փաստաթղթավորել միջերեսները (Աղյուսակ~\ref{tab:ifaces})։
  \task Կազմակերպել ընդհանուր git ռեպոզիտորիա և ինտեգրման թեստեր։
  \task Իրականացնել ամբողջական ինտեգրում և համատեղ ցուցադրում (բաժին~\ref{ch:demo})։
\end{tasklist}

\begin{interfaces}[Միջերեսները երկու մասերի միջև.\label{tab:ifaces}]
  \iface{Տվյալների հոսք}{PS $\rightarrow$ սերվեր}{JSON (բաժին~\ref{ch:data}), \texttt{POST /api/v1/data}, Ethernet}
  \iface{Կոնֆիգուրացիա}{սերվեր $\rightarrow$ PS}{JSON հրաման՝ ռեգիստրի անուն և արժեք (Աղյուսակ~\ref{tab:regmap})}
  \iface{Փորձի կառավարում}{սերվեր $\rightarrow$ PS}{\texttt{start} / \texttt{stop} հրամաններ, \texttt{experiment\_id}}
\end{interfaces}
